\documentclass[11pt]{article}
\usepackage[margin=1in]{geometry}
\usepackage{amsmath,amssymb,amsthm,mathtools}
\usepackage{enumitem}
\usepackage[hidelinks]{hyperref}
\newtheorem{theorem}{Theorem}
\newtheorem{lemma}{Lemma}
\newtheorem{corollary}{Corollary}
\newcommand{\sqcupdot}{\mathbin{\mathaccent\cdot\cup}}
\title{Six-Cycle Insertion and Hook-Unimodality in Higher Lie Characters}
\author{Shawn Calvin Snelling}
\date{22 September 2026}
\begin{document}
\maketitle

\begin{abstract}
Let $M_\lambda(x)$ denote the hook-multiplicity polynomial of the higher Lie character attached to a conjugacy class of cycle type $\lambda$. Adin--Heged\H{u}s--Roichman conjectured that the coefficient sequence of $M_\lambda(x)$ is unimodal for every partition $\lambda$. Using their product formula for rectangular cycle types, we prove the conjecture for the infinite family $(6^s)$ for every $s\ge1$. In fact, $(1+x)M_{(6^s)}(x)$ is strictly log-concave on its nonzero support. We then prove a six-cycle insertion theorem: adjoining any positive number of six-cycles preserves unimodality when the remaining cycle type is unimodal and contains no six. As a consequence, every nonempty partition all of whose parts lie in $\{1,2,3,4,5,6\}$ has a unimodal hook-multiplicity sequence, with no bound on multiplicities or total size.
\end{abstract}

\section{Background and statement of the results}
For a partition $\lambda\vdash n$, write
\[
M_\lambda(x)=\sum_{k=0}^{n-1}m_{(n-k,1^k),\lambda}x^k,
\]
where $m_{(n-k,1^k),\lambda}$ is the multiplicity of the hook character $\chi^{(n-k,1^k)}$ in the higher Lie character $\psi^\lambda$. Adin--Heged\H{u}s--Roichman \cite{AHR} conjectured that the coefficient sequence of $M_\lambda(x)$ is unimodal for every $\lambda$ (their Conjecture 8.1). They verified the conjecture for all $|\lambda|\le15$ and for rectangular shapes $(r^s)$ with $r\le40$ and $s\le5$.

We use two identities from \cite{AHR}. First, when nonempty partitions $\lambda$ and $\nu$ have disjoint sets of part sizes,
\begin{equation}
M_{\lambda\sqcup\nu}(x)=(1+x)M_\lambda(x)M_\nu(x).
\label{eq:disjoint}
\end{equation}
Second, for fixed $r$,
\begin{equation}
1+(1+x)\sum_{s\ge1}M_{(r^s)}(x)y^s
=\prod_{j=0}^r\bigl(1-(-x)^j y\bigr)^{(-1)^{j+1}f_j(r)},
\label{eq:AHRproduct}
\end{equation}
where
\[
f_j(r)=\frac1r\sum_{d\mid\gcd(r,j)}\mu(d)(-1)^{(d+1)j/d}\binom{r/d}{j/d}.
\]
For $r=6$,
\[
(f_0(6),\ldots,f_6(6))=(0,1,3,3,2,1,0),
\]
so
\begin{equation}
E_6(x,y):=1+(1+x)\sum_{s\ge1}M_{(6^s)}(x)y^s
=\frac{(1+xy)(1+x^3y)^3(1+x^5y)}{(1-x^2y)^3(1-x^4y)^2}.
\label{eq:E6}
\end{equation}

\begin{theorem}[all six-cycles]
For every integer $s\ge1$, $M_{(6^s)}(x)$ has unimodal coefficients. Moreover, $(1+x)M_{(6^s)}(x)$ has no internal zeros and is strictly log-concave between its first and last nonzero terms.
\end{theorem}

\begin{theorem}[six-cycle insertion]
Let $\nu$ be a nonempty partition with no part equal to $6$. If $M_\nu(x)$ is unimodal, then for every $s\ge1$, $M_{\nu\sqcup(6^s)}(x)$ is unimodal.
\end{theorem}

\begin{corollary}
Every nonempty partition whose parts belong to $\{1,2,3,4,5,6\}$, at arbitrary multiplicities and arbitrary total size, has a unimodal hook-multiplicity sequence.
\end{corollary}

\section{An exact coefficient formula for $(6^s)$}
Write
\[
M_{(6^s)}(x)=x^{2s-1}P_s(x),\qquad P_s(x)=\sum_{k=0}^{2s+1}q_{s,k}x^k.
\]
The coefficients are
\begin{align}
q_{s,0}&=\frac{s(s+1)}2,\\
q_{s,1}&=\frac{3s^2-s+2}{2},\\
q_{s,2}&=\frac{5s^2-7s+4}{2},
\end{align}
and for $3\le k\le2s$,
\begin{equation}
q_{s,k}=Q_s(k):=
\frac{k^3-4k^2s-3k^2+4ks^2+8ks+6k-4s^2-6s-4}{2},
\label{eq:Qs}
\end{equation}
while
\begin{equation}
q_{s,2s+1}=s.
\end{equation}

For completeness, an exact rational identity yielding these formulas is
\begin{equation}
P_s(x)=\frac{N_0(s,x)+x^{2s}N_1(s,x)}{(1-x)^4},
\label{eq:rationalPs}
\end{equation}
where
\begin{align*}
a_2(x)&=\tfrac12(1-x-x^2+2x^3-x^4-x^5+x^6),\\
a_1(x)&=\tfrac12(1-5x+3x^2-3x^4+5x^5-x^6),\\
a_0(x)&=x-2x^2+5x^3-2x^4+x^5,\\
b_1(x)&=x^5-x^4+x^2-x,\\
b_0(x)&=-2x^4+x^3-2x^2,
\end{align*}
with $N_0(s,x)=a_2(x)s^2+a_1(x)s+a_0(x)$ and $N_1(s,x)=b_1(x)s+b_0(x)$. Summing \eqref{eq:rationalPs} over $s\ge1$ after setting $z=x^2y$ and using the standard sums for $\sum z^s$, $\sum sz^s$, and $\sum s^2z^s$ reproduces \eqref{eq:E6}. Coefficient extraction from $(1-x)^{-4}$ gives the displayed $q_{s,k}$ formulas.

\section{Unimodality of the six-cycle family}
For $s=1,2$ the nonzero coefficient sequences of $P_s$ are respectively
\[
(1,2,1,1),\qquad (3,6,5,5,4,2),
\]
which are unimodal.

Assume $s\ge3$. The first three forward differences are
\[
s^2-s+1,\qquad s^2-3s+1,\qquad \frac{(3s-5)(s-2)}2,
\]
all positive. On the cubic interval, using \eqref{eq:Qs},
\begin{equation}
D_s(k):=Q_s(k+1)-Q_s(k)
=\frac{3k^2-8ks-3k+4s^2+4s+4}{2}
\label{eq:Ds}
\end{equation}
for $3\le k\le2s-1$. This is a convex quadratic in $k$, and
\[
D_s(2s-1)=5-3s<0.
\]
Therefore its nonpositive sublevel set on the interval is contiguous and contains the final point: once a nonpositive forward difference appears, no later forward difference can be positive. Finally,
\[
q_{s,2s+1}-q_{s,2s}=2-2s<0.
\]
Hence the coefficient sequence increases and then decreases, with a possible plateau. Thus $P_s$, and hence $M_{(6^s)}$, is unimodal.

\section{Strict log-concavity after one factor $(1+x)$}
Let
\[
B_s(x)=(1+x)P_s(x)=\sum_{k=0}^{2s+2}b_{s,k}x^k,
\]
and put
\[
L_{s,k}=b_{s,k}^2-b_{s,k-1}b_{s,k+1}.
\]
For $s=1$, $B_1$ has coefficients $(1,3,3,2,1)$ and the internal $L_{s,k}$ are $(6,3,1)$. For $s=2$, the coefficients are $(3,9,11,10,9,6,2)$ and the corresponding differences are $(48,31,1,21,18)$.

Now let $s\ge3$. The first coefficients are
\[
b_0=\frac{s(s+1)}2,\quad b_1=2s^2+1,\quad b_2=4s^2-4s+3,
\quad b_3=\frac{13s^2-25s+18}{2}.
\]
For $4\le k\le2s$,
\begin{equation}
b_k=R_s(k):=
\frac{2k^3-8k^2s-9k^2+8ks^2+24ks+21k-12s^2-24s-18}{2}.
\label{eq:Rs}
\end{equation}
The last two coefficients are $b_{2s+1}=4s-2$ and $b_{2s+2}=s$.

At the first four positions set $u=s-3\ge0$. Direct expansion gives
\begin{align*}
L_{s,1}&=2u^4+24u^3+\tfrac{225}2u^2+\tfrac{483}2u+199,\\
L_{s,2}&=3u^4+29u^3+\tfrac{229}2u^2+\tfrac{433}2u+159,\\
L_{s,3}&=\tfrac94u^4+\tfrac{33}2u^3+\tfrac{233}4u^2+106u+63,\\
L_{s,4}&=9u^4+61u^3+207u^2+335u+181.
\end{align*}
All are positive. For $5\le k\le2s-1$, set $u=k-5\ge0$ and $v=2s-k-1\ge0$. Expansion of \eqref{eq:Rs} gives
\begin{align*}
L_{s,k}={}&2u^2v^2+10u^2v+7u^2+14uv^2+67uv+\tfrac{83}2u\\
&+v^4+10v^3+\tfrac{123}2v^2+172v+118>0.
\end{align*}
At the final two positions, again with $u=s-3$,
\[
L_{s,2s}=13u^2+50u+64>0,
\qquad
L_{s,2s+1}=7u^2+35u+46>0.
\]
Thus $B_s$ is strictly log-concave on its nonzero support. Its endpoint coefficients are positive, so it has no internal zeros.

\section{A convolution lemma and six-cycle insertion}
\begin{lemma}
Let $c$ be a finite nonnegative log-concave sequence with no internal zeros and let $a$ be a finite nonnegative unimodal sequence. Then the convolution $c*a$ is unimodal.
\end{lemma}
\begin{proof}
Extend both sequences by zero to all integers. Log-concavity and contiguous support imply the total-positivity inequality
\[
c_{\ell-i}c_{k-j}\le c_{k-i}c_{\ell-j}\qquad(i<j,\ k<\ell).
\]
Indeed, when the left side is positive this follows by multiplying the nonincreasing successive ratios of $c$; when it is zero the inequality is immediate.

Let $d_j=a_j-a_{j-1}$. Since $a$ is unimodal, all positive terms of $d$ precede all negative terms. Put
\[
P_k=\sum_i d_i^+c_{k-i},\qquad N_k=\sum_j d_j^-c_{k-j}.
\]
The forward difference of $c*a$ at position $k$ is $P_k-N_k$. Multiplying the displayed total-positivity inequality by $d_i^+d_j^-$ and summing gives
\[
P_\ell N_k\le P_kN_\ell\qquad(k<\ell).
\]
A negative forward difference at $k$ followed by a positive forward difference at $\ell$ would require $N_k>P_k$ and $P_\ell>N_\ell$, giving the opposite strict product inequality. Hence the forward differences change sign at most once, and $c*a$ is unimodal.
\end{proof}

Now let $\nu$ be nonempty, unimodal, and contain no part $6$. By \eqref{eq:disjoint},
\[
M_{\nu\sqcup(6^s)}(x)=\bigl((1+x)M_{(6^s)}(x)\bigr)M_\nu(x).
\]
The first factor has a strictly log-concave coefficient sequence with no internal zeros, while the second is nonnegative and unimodal. The lemma proves the six-cycle insertion theorem.

\section{Partitions with parts at most six}
For positive multiplicities, \eqref{eq:AHRproduct} gives
\[
M_{(1^a)}=1,\qquad M_{(2^b)}=x^{2b-1},\qquad M_{(3^c)}=x^{2c-1},
\]
\[
M_{(4^d)}=d\,x^{2d-1}(1+x),
\]
and
\begin{equation}
M_{(5^t)}=x^{2t-1}\left[t+\sum_{j=1}^{2t-1}(2t-j)x^j+x^{2t}\right].
\label{eq:M5}
\end{equation}
The nonzero coefficient sequence in the bracket of \eqref{eq:M5} is
\[
t,\ 2t-1,\ 2t-2,\ldots,1,1,
\]
which is unimodal. Repeated use of \eqref{eq:disjoint} combines distinct part sizes through multiplication by $1+x$, whose coefficient sequence is log-concave. The convolution lemma therefore proves unimodality for every partition with parts at most five. If the partition also contains sixes, remove all six-parts. If a nonempty remainder remains, apply the six-cycle insertion theorem; if no remainder remains, apply the all-six theorem. This proves the corollary.

\section{Scope}
The results above prove an infinite subfamily of Conjecture 8.1 of \cite{AHR}; they do not prove the conjecture for arbitrary partitions. They also do not claim historical priority beyond what can be established by a complete literature review.

\begin{thebibliography}{9}
\bibitem{AHR}
R. M. Adin, P. Heged\H{u}s, and Y. Roichman,
\emph{Higher Lie characters and cyclic descent extension on conjugacy classes},
Algebraic Combinatorics \textbf{6} (2023), no. 6, 1557--1591.
DOI: 10.5802/alco.323.
\end{thebibliography}
\end{document}
